\section{项目概述}

\subsection{项目背景}
在文献更新加速与科研决策压力增大的场景下,核心问题不只是"能不能回答",而是\textbf{可追溯性、可复核性与决策可用性}是否同时成立。
传统关键词检索主要依赖字面匹配,难以覆盖跨语言、跨术语和多步骤研究问题;直接使用通用大模型问答又存在\textbf{幻觉、无出处、不可复现}等风险。
SciScope 的定位不是更强的问答模型,而是将文献资产、检索证据、趋势解释与工具编排形成一条可复用、可验证的科研智能链路,在\textbf{每个结论都附可追溯证据}的前提下支撑文献调研、情报监测与科研治理。

\subsection{应用行业}
高校与科研院所、企业研发部门、科技情报与图书馆机构、科研管理与基金评审部门。
典型场景:研究热点定位、领域脉络梳理、潜在合作者识别、文献调研提速、立项查新与综述辅助。

\subsection{主要工作}
\begin{itemize}
\item \textbf{数据底座}:从 arXiv、PubMed、PMC、OpenAlex、Crossref、DOAJ 6 个公开来源采集、清洗、规范化、去重,形成 168{,}447 条分析语料、159{,}164 篇处理后语料文件和 159{,}187 篇运行时入库论文,覆盖计算机/生物医学/材料等多领域、主窗口 2022--2026。
\item \textbf{数据分析报告}(交付物 1):文献分布、关键词演化、作者合作网络、主题模型、趋势统计(见同级交付物 \texttt{sciscope\_data\_report})。
\item \textbf{科研智能体模型}(交付物 2):本地嵌入索引、混合检索、趋势预测、论文推荐、知识图谱五大模型,配套 FastAPI 服务与终端智能体客户端。
\item \textbf{系统工程}:PostgreSQL + pgvector 服务层、本地大模型(vLLM/MLX)生成层、一键复现的 Makefile 流水线与自动化测试。
\end{itemize}

\subsection{创新点}
\begin{reportbox}{五个核心创新}
\begin{enumerate}
\item \textbf{大模型无关的证据接地架构(LLM-agnostic, evidence-grounded)}:大语言模型(LLM)只负责"用人话写答案",领域智能沉淀在本地可复现的索引/模型文件中;生成层经 OpenAI 兼容接口\textbf{可插拔}(本地 7B ↔ DeepSeek 云端,一处配置切换),换任意大模型系统照常工作。
\item \textbf{工具契约与确定性校验门}:9 个原生语料工具与 1 个 \texttt{delegate} 子智能体工具(共 10 个运行时工具)均通过统一 \texttt{Tool} 注册表暴露;工具声明只读性、参数校验和结果上限。调用前可拦截模型编造的论文主键(\texttt{paper\_id});显式科研技能还设置首工具强制调用与工具预算,从架构上降低幻觉、空跑和捏造风险,便于审计。
\item \textbf{字面 + 语义混合检索(Hybrid Retrieval)}:PostgreSQL 全文检索(FTS)与 pgvector 向量检索双路并行,用倒数排名融合(Reciprocal Rank Fusion, RRF)合并,兼顾精确匹配与语义/跨语言理解(中文问句可命中英文文献)。
\item \textbf{"模型文件"即可复现资产}:不依赖黑盒训练权重,而是用本项目代码从本项目语料\textbf{计算/拟合}出的嵌入索引、趋势模型、推荐向量、知识图谱;\texttt{make agent-build} 可一键重建,知识产权清晰、评委可复现。
\item \textbf{MCP 双向互通与专员子智能体}:SciScope 可作为模型上下文协议(Model Context Protocol, MCP)服务端把工具开放给外部客户端,也可消费外部 MCP 工具;复杂任务可通过 \texttt{delegate} 下派给综述员、趋势分析师、批判核查员三类专员子智能体,且结构上禁止递归委派。
\item \textbf{可观测运行与数据质量工程}:LangGraph 节点事件携带阶段、耗时、终止原因(\texttt{stop\_reason})和估算 token 用量;新增 \texttt{make agent-smoke} 黑盒验收,持续检查真实库表规模、能力边界、论断核查、趋势分析和论文推荐的工具链路。数据层通过标题优先去重、front-matter 剔除、PDF 垃圾过滤和关键词噪声过滤提升检索、趋势与图谱信噪比。
\end{enumerate}
\end{reportbox}

\begin{plainbox}
\textbf{数字口径说明}:本报告将数据规模分为四类口径:数据分析报告使用 \texttt{data/analysis/summary.json} 的 168{,}447 条分析语料;处理后语料文件使用 \texttt{data/processed/papers\_corpus.summary.json} 的 159{,}164 篇;运行时服务以 PostgreSQL 当前库表为准,即 159{,}187 篇论文、367{,}773 个片段、367{,}773 条片段向量和 159{,}135 条论文级向量;系统效果以 \texttt{output/eval/eval\_report.json} 为准。四类口径用途不同,不混写为同一个总数。
\end{plainbox}
